\chapter*{Abstract}
Privacy and publicity are commonly treated as opposing values: privacy conceals, whereas publicity exposes. This essay argues that both are responses to the same underlying condition---the severe difficulty of reliable judgment in causally dense environments. Public recording can preserve enormous quantities of evidence without producing timely understanding; indeed, the accumulation of video, testimony, metadata, and commentary may enlarge the interpretive search space faster than investigators can reduce it. Salient facts may take years or decades to move through hierarchies of attention and acquire legal significance. Under these conditions, suspicion is neither simply irrational nor sufficient for action. Its legitimate function is to keep open the possibility that appearances are incomplete. Legal requirements that appear cumbersome---warrants, evidentiary standards, adversarial testing, burdens of proof, and appellate delay---constitute an institutional form of disciplined paranoia. They assume that observers, officials, crowds, and technical systems may be mistaken or untrustworthy. By slowing the passage from suspicion to coercion, these constraints do not merely limit power; they create the protected uncertainty within which freedom becomes possible.
